\section{Ubicación de la clase: la multimodalidad no consiste en añadir un módulo de visión por fuera del LLM}

Esta clase abarca mucho terreno. La primera parte repasa tres «momentos» de los modelos de lenguaje: BERT, GPT-3 y ChatGPT. La parte intermedia trata la arquitectura Transformer, los sistemas de entrenamiento, la alignment y la ingeniería de datos. Solo la última parte llega a BLIP-2, LLaVA, CogVLM, CogAgent, la difusión de imágenes y la generación de video. Este orden no es una digresión, sino la respuesta a una pregunta que recorre toda la clase: \textbf{cuando las capacidades del modelo se extienden de text a image, GUI y video, ¿qué regularidades provenientes de los LLM siguen siendo válidas y cuáles exigen cambiar de representación, de función objetivo o de ruta de sistema?}

\subsection{Corrección de fuentes y límites de la clase}

El borrador anterior fechaba la clase el 4 de abril de 2026 y añadía mucho contenido que no existe en el curso, como monitoring dashboard, drift detection, automatic rollback y fairness governance. Ahora esta clase toma como eje la página oficial del curso CS25 V4, la sesión del 9 de mayo de 2024, la grabación oficial de Stanford Online \texttt{cYfKQ6YG9Qo}, los subtítulos humanos oficiales en \texttt{en-US} y 31 teaching states recuperados de la grabación. Ni la página oficial ni la descripción del video ofrecen un slide deck independiente que pueda verificarse; por eso, la pantalla compartida en la grabación es la source of truth visual.

\lecturefigure{slide-01-title.jpg}{Título oficial de la clase: From Large Language Models to Large Multi-modality Models}{Grabación oficial de Stanford Online 00:02:31}

El término ``Large Multi-modality Models'' del título es el que se usaba en la clase de 2024. No significa que todas las modality ya se hayan resuelto de manera unificada; subraya que la investigación sobre modelos ya pasó de una sola interfaz de lenguaje a la comprensión visual, la generación de imágenes, las GUI action y el video. CogVLM, CogAgent, CogView y CogVideo, del equipo del ponente, son casos de estudio y no la única conclusión del curso.

\begin{warningbox}{Hay que conservar el límite temporal}
Los benchmark, las capacidades de producto, las cifras de descargas y las predicciones «a uno o dos años» que se muestran en esta clase son una instantánea de la sesión del 2024-05-09. Las versiones posteriores de los modelos, los cambios en las tablas de clasificación o los resultados de la industria no deben proyectarse hacia atrás como hechos de la clase; si se usan como ampliación, deben marcarse explícitamente como post-lecture evidence.
\end{warningbox}

\subsection{Hoja de ruta en cuatro partes y preguntas de lectura}

El curso primero pregunta why, luego how, después what y por último where. El lector no debe ver las cuatro partes como una revisión de temas paralelos. El objective y el scaling de los LLM explican por qué la ingeniería de sistemas cobró importancia. Los sistemas y los datos explican por qué los VLM pueden reutilizar con rapidez la base de lenguaje. A su vez, las carencias de la representación visual explican por qué la generación pasó de la autoregression a la diffusion.

\lecturefigure{slide-02-outline.jpg}{Hoja de ruta del curso: historia, técnicas de entrenamiento, modelos de visión y lenguaje y líneas de investigación}{Grabación oficial de Stanford Online 00:03:29}

\begin{importantbox}{Cuatro preguntas para leer toda la clase}
\begin{enumerate}
  \item \textbf{Objective:} ¿las distintas modality deben compartir un autoregressive objective o usar un objective especializado de diffusion / contrastive?
  \item \textbf{Architecture:} ¿qué problema de interfaz resuelven, cada uno, la capa de proyección, el Q-Former, los vision experts y la cross attention?
  \item \textbf{Systems:} ¿cómo cambian el límite de los problemas entrenables ZeRO (Zero Redundancy Optimizer, que reparte los optimizer/gradient/parameter states entre los data-parallel ranks), el parallelism, el long context y los token de alta resolución?
  \item \textbf{Data:} cuando los web data generales se acercan a la saturación, ¿cómo se convierte el compute en datos de mayor calidad y verificables?
\end{enumerate}
\end{importantbox}

\teachervoice{00:02:31--00:04:35. Ming Ding advierte al público que la clase recorrerá varios subcampos que quizá no conozca bien. Su objetivo no es hacer un survey completo de cada modelo, sino explicar, siguiendo la cronología de la comunidad Transformer, por qué se desplazó el foco de la investigación.}

\subsection{Registro de evidencia: mecanismos, resultados, juicios y predicciones}

Una misma slide puede contener a la vez un diagrama de arquitectura de un artículo, un juicio expresado en clase y un resultado de producto. Para evitar que «lo dice la figura» se lea como un hecho general, esta clase usa un registro de evidencia en cuatro niveles.

\begin{knowledgebox}{Cómo leer la evidencia en cuatro niveles}
\begin{tabularx}{\textwidth}{L{0.18\textwidth} L{0.31\textwidth} X}
\toprule
Nivel & Ejemplos & Lectura correcta \\
\midrule
Mecanismo estándar & cross entropy, attention, DDPM, DPO & Dar fórmulas, variables y supuestos \\
Sistema concreto & ZeRO, Megatron, CogVLM, MM-DiT & Explicar interfaces, rutas de recursos y versiones \\
Resultado de clase & benchmark table, web-agent trace, cifras de velocidad & Solo respaldan la configuración y el protocol mostrados \\
Juicio/predicción del ponente & la loss determina las capacidades; el video será un tema central & Conservar la fecha, el tono y el margen para contraejemplos \\
\bottomrule
\end{tabularx}
\end{knowledgebox}

\subsection{Resumen de la sección}

Esta es una clase sobre «trasladar las regularidades de los LLM a la multimodality», no un manual de producto de la serie Cog. El lector debe seguir a la vez el objective, la architecture, los systems y los data, y distinguir siempre entre mecanismos estándar, experimentos de clase, juicios del ponente y predicciones a futuro.
